\section{Estimating Equations and Assumptions}
\label{app:estimating-equations}

This appendix gives the estimating equations for the three causal estimands defined in Section~\ref{sec:estimands}. Appendix~\ref{app:experimental-details} provides the implementation details for each experiment.

\subsection{Notation}

We follow the notation introduced in Section~\ref{sec:methods}, adding an index $i=1,\ldots,n$ for the $n$ completed trajectories.

\paragraph{Trajectories and intervention points.}
For trajectory $i$, let $T_i$ denote the number of intervention points, indexed by $\mathcal T_i=\{1,\ldots,T_i\}$.
At each point $t\in\mathcal T_i$, let $H_{it}$ denote the pre-assignment history defined in~\eqref{eq:history}, and let $A_{it}\in\{0,1\}$ denote the randomized intervention assignment.
We use $Y_i$ for the terminal outcome, which is binary in $\tau$-bench Retail and ImpossibleBench and continuous in ToolSandbox.
For brevity, we write $\sum_{i,t}$ for $\sum_{i=1}^n\sum_{t\in\mathcal T_i}$.

\paragraph{Intervention policies.}
Let $g$ denote the randomized collection policy, with assignment probabilities
\[
g(a\mid h)=e(h)^a(1-e(h))^{1-a},
\qquad
e(h)=\Pr_g(A_t=1\mid H_t=h).
\]
In our experiments, we use $e(h)=1/2$ at every eligible intervention point.
We write $\E_g$ for expectation over trajectories generated under $g$.
We use $\pi$ to denote a target intervention policy, including the never-intervene policy $\pi_0$ and the selective policy $\pi_G$ defined in~\eqref{eq:policy-target}.

\paragraph{Reporting groups and estimands.}
For a reporting group $G$, let $I_{it}(G)=\ind\{H_{it}\in G\}$ indicate whether intervention point $t$ in trajectory $i$ belongs to $G$.
Define the number of matching intervention points by
\[
N_i(G)=\sum_{t\in\mathcal T_i}I_{it}(G),
\qquad
N(G)=\sum_{i=1}^n N_i(G).
\]
We consider three causal estimands: the blip with randomized continuation $\tau(G;g)$, the blip with no-treatment continuation $\tau(G;0)$, both defined in~\eqref{eq:blip}, and the policy gain $\Delta(G)$ defined in~\eqref{eq:policy-target}.

\paragraph{Cross-fitting.}
We divide the $n$ trajectories into five folds, keeping all intervention points from the same trajectory together.
Let $k(i)\in\{0,\ldots,4\}$ denote the fold containing trajectory $i$, and define $\mathcal I_k=\{i:k(i)=k\}$.
We use the superscript $(-k)$ to denote a function fitted without data from fold $k$, so predictions for trajectory $i$ are obtained from models fitted without its fold $k(i)$.

\subsection{Blip with Randomized Continuation}
\label{app:aipw}

\paragraph{Augmented inverse probability weighting.}
Let $m_g(h,a)=\E_g[Y\mid H_t=h,A_t=a]$ denote the conditional outcome regression, and let $\widehat m_g^{(-k)}(h,a)$ be its estimate fitted without fold $k$.
By~\eqref{eq:identification}, the blip effect is $b(h;g)=m_g(h,1)-m_g(h,0)$.
We then define the cross-fitted augmented inverse probability weighting (AIPW) score as
\begin{align}
    \widehat D_{it}
    ={}&\widehat m_g^{(-k(i))}(H_{it},1)-\widehat m_g^{(-k(i))}(H_{it},0) +\frac{A_{it}}{e(H_{it})}\bigl\{Y_i-\widehat m_g^{(-k(i))}(H_{it},1)\bigr\}\nonumber\\
    &-\frac{1-A_{it}}{1-e(H_{it})}\bigl\{Y_i-\widehat m_g^{(-k(i))}(H_{it},0)\bigr\},
    \label{eq:aipw}
\end{align}
which gives the group estimator for group $G$
$$
    \widehat\tau_{\mathrm{AIPW}}(G;g)=\frac{1}{N(G)}\sum_{i,t}I_{it}(G)\widehat D_{it}.
$$


When the propensity $e(h)$ is known, the score stays conditionally unbiased even when the outcome predictions are
wrong. Let $r_a$ be any regression fitted
without the current trajectory. At eligible histories, randomization gives
\[
    \E_g\!\left[\frac{\ind\{A_t=a\}}{g(a\mid H_t)}\{Y-r_a(H_t)\}\,\middle|\,H_t=h\right]
    =m_g(h,a)-r_a(h),
\]
so the right-hand side of~\eqref{eq:aipw} with each $\widehat m_g$ replaced by
$r_a$ has conditional mean $b(h;g)$.  When assignment probabilities are known, the outcome regression affects the
estimator's variance but does not change its population target. In case of estimated propensities the usual double-robustness
property applies \citep{bang2005dr}.

\paragraph{Linear blip model for randomized continuation}
Let $f(h)$ denote an encoding of history $h$ used in the effect model. It includes
attributes such as the proposed tool, earlier errors, and prior interventions.
Within each environment and estimand, all agents use the same feature definition.

Let $s(h)\in\{0,\ldots,P-1\}$ identify the phase and write $s_{it}=s(H_{it})$.
$\tau$-bench Retail and ToolSandbox use semantic progress phases (see Appendix~\ref{app:phases-retail} and Appendix~\ref{app:phases-tool}, correspondingly), while steering uses
chronological checkpoints. The features and phases for each experiment are defined in
Appendix~\ref{app:experimental-details}.

Let $d=\dim f(h)$ and partition the coefficient coordinates into a shared set
$\mathcal S\subseteq\{1,\ldots,d\}$ and a phase-specific set
$\mathcal U=\{1,\ldots,d\}\setminus\mathcal S$. Define
\begin{equation}
 \beta_{s,\mathcal S}=\alpha,\qquad
 \beta_{s,\mathcal U}=\gamma_s,\qquad s=0,\ldots,P-1,
 \label{eq:coefficient-sharing}
\end{equation}
where $\alpha\in\mathbb R^{|\mathcal S|}$ is shared across phases and
$\gamma_s\in\mathbb R^{|\mathcal U|}$ may vary with $s$.  The cases $\mathcal S=\varnothing$ and $\mathcal U=\varnothing$
give fully phase-specific and fully shared coefficients, respectively. The exact choice of these coefficients for our experiments is described in Appendix~\ref{app:experimental-details}.

Define
\[
 q(h)=\E_g[Y\mid H_t=h]=\{1-e(h)\}m_g(h,0)+e(h)m_g(h,1),
 \qquad Z_{it}=(A_{it}-e(H_{it}))f(H_{it}).
\]
The cross-fitted value $\widehat q_{it}:=\widehat q^{(-k(i))}(H_{it})$.
Under a well specified model we have that the blip with randomized continuation $b(h;g)=f(h)^\top\beta_{s(h)}(g)$.

Finally, we estimated the coefficients in our model jointly by minimizing the following objective:
\begin{equation}
\begin{aligned}
 \bigl(\widehat\alpha(g),\{\widehat\gamma_s(g)\}\bigr)
   &=\argmin_{\alpha,\{\gamma_s\}}\sum_{i,t}
       \bigl\{Y_i-\widehat q_{it}-Z_{it}^\top\beta_{s_{it}}(g)\bigr\}^2,\\
 \widehat\tau_{\mathrm{SNMM}}(G;g)
   &=\frac{1}{N(G)}\sum_{i,t}I_{it}(G)f(H_{it})^\top\widehat\beta_{s_{it}}(g).
\end{aligned}
\label{eq:local-snmm}
\end{equation}
This algorithm for estimating the coefficients $\beta$ is known in the literature as the Double Machine Learning algorithm \cite{chernozhukov2018dml} or equivalently the R-Learner algorithm \cite{nie2021quasi} with a linear specification of the heterogeneous treatment effect (see e.g. the `LinearDML' class in the EconML library \cite{econml}). 


In case the linear model is misspecified, the coefficients target an $e(h)\{1-e(h)\}$-weighted projection of the true blip onto the allowed effect model. When the propensities are constant, as it is in our case, then this weighted projection reduces to the un-weighted projection of the true blip on the allowed effect model. This explains why the SNMM estimatein our experiments is almost identical to the AIPW estimate.

\subsection{Blip with No-treatment Continuation}
\label{app:snmm}

\paragraph{Shared and phase-specific effects.}
For each changed continuation $\pi$, assemble $\beta_s(\pi)$ from
$\alpha(\pi)$ and $\gamma_s(\pi)$ using~\eqref{eq:coefficient-sharing}.
The structural restriction is
\begin{equation}
 \E\!\left[Y_t(a;\pi)-Y_t(0;\pi)\mid H_t=h\right]
   =a f(h)^\top\beta_{s(h)}(\pi),\qquad a\in\{0,1\}.
 \label{eq:structural}
\end{equation}
Here $Y_t(a;\pi)$ retains the observed past, assigns $a$ now, and follows $\pi$
afterward. We fit a separate model for each target continuation, sharing the
coefficients in $\mathcal S$ across phases and allowing those in $\mathcal U$
to vary by phase. We adapt the SNMM g-estimation construction of
\citet{lewis2021dynamic}; see also
\citet{vansteelandt2014snmm}.

\paragraph{Future-treatment correction.} Keep the centered current features $Z_{it}=(A_{it}-e(H_{it}))f(H_{it})$.
For each phase $\ell$, collect the features of later decisions in that phase:
\begin{equation}
 B_{it,\ell}(\pi)
   =\sum_{\substack{j\in\mathcal T_i\\j>t,\ s_{ij}=\ell}}
       \{A_{ij}-\pi(H_{ij})\}f(H_{ij}).
 \label{eq:future}
\end{equation}
 Define
$p_\ell(h,\pi)=\E_g[B_{t,\ell}(\pi)\mid H_t=h]$ and the cross-fitted residual
\[
 \widehat U_{it,\ell}(\pi)
   =B_{it,\ell}(\pi)-\widehat p_\ell^{(-k(i))}(H_{it},\pi).
\]
All these nuisance regressions condition on the history at the current decision $H_{it}$. 

\paragraph{Joint coefficient estimation.}
For candidate coefficients $\alpha$ and $\{\gamma_s\}$, let
\begin{equation}
 \varepsilon_{it}(\alpha,\{\gamma_s\};\pi)
   =Y_i-\widehat q_{it}-Z_{it}^\top\beta_{s_{it}}(\pi)
     -\sum_{\ell=0}^{P-1}\widehat U_{it,\ell}(\pi)^\top\beta_\ell(\pi),
 \label{eq:joint-snmm-residual}
\end{equation}
where each $\beta_\ell(\pi)$ is assembled from these candidate coefficients.
Write $\widehat\varepsilon_{it}(\pi)$ for this residual evaluated at
$\widehat\alpha(\pi)$ and $\{\widehat\gamma_s(\pi)\}$. We estimate them jointly
by solving
\begin{equation}
\begin{aligned}
 \sum_{i,t}Z_{it,\mathcal S}\widehat\varepsilon_{it}(\pi)&=0,\\
 \sum_{i,t:s_{it}=s}Z_{it,\mathcal U}\widehat\varepsilon_{it}(\pi)&=0,
    \qquad s=0,\ldots,P-1.
\end{aligned}
\label{eq:joint-snmm-moments}
\end{equation}
Thus each shared coefficient pools its score across phases, while each
phase-specific coefficient uses its own phase's score. 

When $\mathcal S=\varnothing$, coefficients are shared across decisions
within each phase but remain separate across phases. If phases are
nondecreasing along a trajectory, later decisions cannot belong to an
earlier phase. We enforce this restriction by setting
$\widehat U_{it,\ell}(\pi)=0$ for $\ell<s_{it}$. The joint system
in~\eqref{eq:joint-snmm-moments} is then block upper triangular and can
be solved by backward substitution, one phase at a time. When the
diagonal blocks are nonsingular, this gives the same coefficients as
solving the full system at once.

We use this backward implementation for $\tau$-bench Retail
(Appendix~\ref{app:phases-retail}) and for the separate chronological
checkpoints in ImpossibleBench (Appendix~\ref{app:steering-fit}).
ToolSandbox shares tool and history coefficients across phases and
instead solve a joint system of linear equations (Appendix~\ref{app:phases-tool}).
The backward implementation adapts Algorithm~3 of
\citet{lewis2021dynamic} to allow repeated decisions within a phase.

For $s=P-1,\ldots,0$, we first subtract the contributions of later
phases using their already estimated coefficients:
\begin{equation}
 \overline Y_{it,s}(\pi)
   =Y_i-\widehat q_{it}
     -\sum_{\ell>s}\widehat U_{it,\ell}(\pi)^\top
       \widehat\beta_\ell(\pi),
 \qquad s_{it}=s.
 \label{eq:backward-outcome}
\end{equation}
We then solve the current phase's linear moment equations:
\begin{equation}
 \sum_{i,t:s_{it}=s}Z_{it}
  \Bigl[\overline Y_{it,s}(\pi)
    -\{Z_{it}+\widehat U_{it,s}(\pi)\}^\top
      \widehat\beta_s(\pi)\Bigr]=0.
 \label{eq:phase-moments}
\end{equation}
The additional term $\widehat U_{it,s}(\pi)$ accounts for later
decisions that remain in the current phase. For example, a Retail
episode may contain several proposals before successful customer
identification. These decisions share the same coefficient vector
$\beta_s(\pi)$, so their future contributions must remain in the
current phase's equation while that vector is estimated. We therefore
solve jointly over all decisions in phase $s$, retaining both the
current term $Z_{it}$ and the future term
$\widehat U_{it,s}(\pi)$. In ImpossibleBench, each chronological checkpoint occurs at most once, so we set this same-phase future term to zero.

When the treatment propensity $e(h)$ is known, the above estimating equations recover the true SNMM parameters even if the future-treatment regressions $\widehat p_\ell$ and the outcome regressions $\widehat q_{it}$ are misspecified. This is because $\E[r(H_{it})Z_{it}]=0$, for any function $r$. Hence, in expectation these centering regression terms vanish. These estimated regressions are used for variance reduction.



\paragraph{Identification.}
Let $\nu(h,\pi)=\E[Y_t(0;\pi)\mid H_t=h]$. Under the assumption that the SNMM model is well specified, removing
the later treatment effects gives
\begin{equation}
 \E_g\!\left[Y-\sum_{\ell=0}^{P-1}B_{t,\ell}(\pi)^\top\beta_\ell(\pi)
       \,\middle|\,H_t=h,A_t=a\right]
   =\nu(h,\pi)+a f(h)^\top\beta_{s(h)}(\pi).
 \label{eq:snmm-recursion}
\end{equation}
Subtracting the current centered effect $Z_t^\top\beta_{s(H_t)}(\pi)$ from the
corrected outcome and multiplying by $Z_t$ gives zero expected phase scores at
the true coefficients. Pooling their shared coordinates across phases gives the population
version of~\eqref{eq:joint-snmm-moments}. Since $\E_g[Z_t\mid H_t]=0$, subtracting
pre-assignment controls for $q$ and $p_\ell$ leaves these moments unchanged under
known randomization; the controls affect precision.

\paragraph{Blip with no-treatment continuation.}
After fitting with $\pi=\pi_0$, average the fitted effects over matching
opportunities:
\begin{equation}
 \widehat\tau(G;0)=\frac{1}{N(G)}\sum_{i,t}I_{it}(G)
     f(H_{it})^\top\widehat\beta_{s_{it}}(\pi_0).
 \label{eq:no-future-estimate}
\end{equation}
This retains the histories reached under the collection policy. Each selective
policy $\pi_G$ receives a separate fit with the same features, phases, and sharing
pattern; its gain uses the fitted $\pi_0$ reference. The primary comparison for the blip with randomized
continuation remains the local fit~\eqref{eq:local-snmm}, without future correction.

\subsection{Policy Value Estimation}
\label{app:value}

The coefficients fitted for continuation $\pi$ also identify its value:
\begin{equation}
 V(\pi)=\E_g\!\left[Y+\sum_{t\in\mathcal T}
    \{\pi(H_t)-A_t\}f(H_t)^\top\beta_{s(H_t)}(\pi)\right].
 \label{eq:value-identity}
\end{equation}
Each summand removes the effect of the observed assignment and adds that of the target, using the coefficients for the phase at that history and the target continuation \citep[Lemma~6]{lewis2021dynamic}. The estimator and policy gain are
\begin{equation}
\begin{aligned}
 \widehat V(\pi)&=\frac1n\sum_{i=1}^n\left[Y_i+
    \sum_{t\in\mathcal T_i}\{\pi(H_{it})-A_{it}\}
      f(H_{it})^\top\widehat\beta_{s_{it}}(\pi)\right],\\
 \widehat\Delta(G)&=\widehat V(\pi_G)-\widehat V(\pi_0).
\end{aligned}
\label{eq:value-estimator}
\end{equation}

